\documentclass[11pt,a4paper]{article}
\usepackage{shared/luxcover}
\usepackage[utf8]{inputenc}
\usepackage[T1]{fontenc}
\usepackage{amsmath,amssymb,amsfonts,amsthm}
\usepackage{graphicx}
\usepackage{hyperref}
\usepackage{booktabs}
\usepackage{algorithm}
\usepackage{algpseudocode}
\usepackage{listings}
\usepackage{xcolor}
\usepackage{tikz}
\usepackage{float}
\usepackage{enumitem}
\usepackage{array}
\usepackage{multirow}
\usepackage{tabularx}
\input{shared/lstlang}
\usetikzlibrary{arrows.meta,positioning,shapes.geometric,fit,calc}

\hypersetup{colorlinks=true,linkcolor=black,citecolor=blue,urlcolor=blue}

\lstset{
  basicstyle=\ttfamily\scriptsize,
  breaklines=true,
  frame=single,
  numbers=left,
  numberstyle=\tiny\color{gray},
  keywordstyle=\color{blue!70!black}\bfseries,
  commentstyle=\color{green!50!black}\itshape,
  stringstyle=\color{red!60!black},
  showstringspaces=false,
}

\newtheorem{theorem}{Theorem}
\newtheorem{lemma}[theorem]{Lemma}
\newtheorem{definition}{Definition}
\newtheorem{proposition}[theorem]{Proposition}
\newtheorem{remark}{Remark}

\title{Non-Custodial Blockchain ATS: Architecture of a\\Regulatory-Compliant Securities Exchange}

\author{
Zach Kelling\thanks{Corresponding author: zach@lux.network}\\
\textit{Lux Partners}\\
\textit{A subsidiary of Hanzo Industries Inc}\\
\texttt{research@lux.network}
}

\date{2025}

\begin{document}

\luxcoverpage

\begin{abstract}
We present the architecture of a non-custodial Alternative Trading System (ATS) built on Lux Network, designed to trade digital securities in full compliance with SEC Regulation ATS, broker-dealer requirements under the Securities Exchange Act of 1934, and transfer agent obligations under Section 17A. The system is implemented as a single Go binary running on Hanzo Base, comprising 32 collections, 64 API routes, and 19 scheduled tasks. Non-custodial operation is achieved through MPC threshold signing (\texttt{luxfi/mpc}) with a 2-of-3 threshold requirement for any asset transfer---the platform never holds private keys. This work descends directly from Arca's ArCoin infrastructure (2018), which pioneered interest-bearing digital securities backed by US Treasury bills using ERC-1400/ERC-1404 security token contracts. We describe the complete system: order matching, settlement on Lux EVM with T+0 atomic finality, KYC/AML integration via Partner, equities routing through Alpaca OmniSub, and regulatory reporting. We compare the architecture against existing platforms (tZERO, Securitize Markets, INX) and demonstrate that non-custodial MPC custody eliminates the single largest attack surface in securities infrastructure while preserving regulatory compliance.
\end{abstract}

\section{Introduction}

The United States securities market operates under a regulatory framework designed for centralized intermediaries. Broker-dealers hold customer securities in street name. Transfer agents maintain issuer registries on private databases. Alternative Trading Systems match orders through proprietary systems that custody assets during settlement. Every intermediary holds keys, creating attack surfaces that have produced losses exceeding \$10B in the traditional finance sector over the past decade through insider theft, operational failures, and cyberattacks.

The blockchain industry has attempted to address these structural risks by building ``decentralized exchanges,'' but DEX architectures designed for fungible utility tokens are fundamentally incompatible with the regulatory requirements of securities trading. Securities require transfer restrictions, accreditation verification, jurisdictional compliance, forced transfers for regulatory actions, and auditable cap table management. No existing DEX satisfies these requirements.

\subsection{From Arca to Lux}

This work originates in the Arca platform (ar.ca), where the author served as CTO in 2018. Arca built several foundational components for digital securities infrastructure:

\begin{itemize}[nosep]
  \item \textbf{ArCoin}: An interest-bearing digital security backed by US Treasury bills, implemented as an ERC-1400 token with automated NAV calculation and distribution.
  \item \textbf{st-contracts}: A security token protocol implementing ERC-1400 (Security Token Standard) and ERC-1404 (Simple Restricted Token Standard) with on-chain transfer restrictions.
  \item \textbf{exchange-sdk}: A trading SDK for digital securities with order book management, matching engine, and settlement coordination.
  \item \textbf{sec-wallet}: An institutional-grade wallet for securities custody with multi-signature authorization and audit logging.
  \item \textbf{Arca Funds}: The first crypto hedge fund structure with tokenized fund interests represented as security tokens.
  \item \textbf{miners/binance-recorder}: Market data infrastructure for recording and replaying exchange data feeds.
\end{itemize}

The central lesson from Arca was that custodial architectures---even well-engineered ones---create irreducible counterparty risk. When a platform holds private keys on behalf of investors, every operational failure, insider threat, and regulatory seizure becomes a potential loss event. The path forward required eliminating custodial risk at the protocol level.

This paper describes that path: a non-custodial ATS built on Lux Network that satisfies every requirement of SEC Regulation ATS, broker-dealer registration, and transfer agent registration while never holding a single private key.

\subsection{Contributions}

\begin{enumerate}[nosep]
  \item A complete ATS architecture that is non-custodial by construction, using MPC threshold signatures for all asset movements.
  \item A single-binary deployment model on Hanzo Base that packages order matching, settlement, compliance, and reporting into 32 collections and 64 routes.
  \item Atomic T+0 settlement on Lux EVM, eliminating the T+1/T+2 settlement risk present in traditional securities infrastructure.
  \item A regulatory compliance framework that maps SEC/FINRA requirements to on-chain enforcement mechanisms.
  \item Integration architecture for hybrid traditional/digital securities via Alpaca OmniSub.
\end{enumerate}

\subsection{Related Work}

The Lux MPC custody architecture is detailed in the companion paper on M-Chain \cite{lux-mchain-mpc}. The security token standard used for on-chain securities is formalized in the Lux Security Token Standard \cite{lux-security-token-standard}. Smart order routing across multiple execution venues is described in the Lux Smart Order Routing paper \cite{lux-smart-order-routing}. The transfer agent subsystem is presented separately \cite{lux-transfer-agent}. The Lux credit protocol, which enables margin lending against tokenized securities, is specified in \cite{lux-credit-protocol-spec}.

\section{Regulatory Framework}\label{sec:regulatory}

\subsection{ATS Registration (Regulation ATS)}

An Alternative Trading System under SEC Rule 300(a) is any organization that constitutes, maintains, or provides a marketplace for bringing together purchasers and sellers of securities. The Lux ATS registers on Form ATS-N with the SEC and operates under the following obligations:

\begin{enumerate}[nosep]
  \item \textbf{Fair access}: Any qualified investor meeting accreditation and KYC/AML requirements may trade.
  \item \textbf{Order display}: Orders above applicable thresholds are displayed to all subscribers.
  \item \textbf{Best execution}: The smart order router satisfies Regulation NMS best execution obligations \cite{lux-smart-order-routing}.
  \item \textbf{Record keeping}: All orders, executions, and cancellations are logged immutably on-chain.
  \item \textbf{System integrity}: The platform maintains business continuity and disaster recovery per Regulation SCI.
\end{enumerate}

\subsection{Broker-Dealer Registration}

The operating entity registers as a broker-dealer under Section 15 of the Securities Exchange Act of 1934. This registration carries obligations for:

\begin{itemize}[nosep]
  \item Net capital requirements (Rule 15c3-1)
  \item Customer protection (Rule 15c3-3)
  \item Anti-money laundering compliance (Bank Secrecy Act, USA PATRIOT Act)
  \item Books and records (Rules 17a-3, 17a-4)
  \item Supervision (FINRA Rules 3110, 3120)
\end{itemize}

The non-custodial architecture satisfies Rule 15c3-3 in a novel manner: because the platform never holds customer securities or funds (all assets remain in MPC-custodied wallets controlled by the investor), the customer protection rule's possession-or-control requirements are satisfied by the investor's own cryptographic control.

\subsection{Transfer Agent Registration}

The platform operates as a registered transfer agent under Section 17A of the Securities Exchange Act. Transfer agent obligations and their on-chain implementations are described in the companion paper \cite{lux-transfer-agent}.

\section{System Architecture}\label{sec:architecture}

\subsection{Deployment Model}

The ATS is implemented as a single Go binary (\texttt{luxfi/ats}) running on Hanzo Base. Hanzo Base provides an embedded SQLite-compatible database with real-time subscriptions, authentication, and file storage. The binary packages all ATS functionality into a single deployable unit.

\begin{table}[H]
\centering
\caption{ATS system components.}
\label{tab:components}
\begin{tabular}{lrl}
\toprule
\textbf{Component} & \textbf{Count} & \textbf{Description} \\
\midrule
Collections & 32 & Data models (orders, securities, accounts, \ldots) \\
API routes & 64 & REST endpoints (/v1/\ldots) \\
Cron jobs & 19 & Scheduled tasks (settlement, compliance, \ldots) \\
Event hooks & 47 & Real-time triggers on collection changes \\
MPC signers & 3 & Threshold signature participants \\
\bottomrule
\end{tabular}
\end{table}

\subsection{Collection Schema}

The 32 collections partition the system state into distinct concerns:

\begin{table}[H]
\centering
\caption{Collection groups and their members.}
\label{tab:collections}
\begin{tabular}{lp{8cm}}
\toprule
\textbf{Group} & \textbf{Collections} \\
\midrule
Identity & accounts, investors, accreditations, kyc\_checks, aml\_screens \\
Securities & securities, offerings, partitions, documents, corporate\_actions \\
Trading & orders, executions, order\_books, quotes, rfqs \\
Settlement & settlements, transfers, mpc\_requests, mpc\_signatures \\
Compliance & compliance\_events, regulatory\_reports, audit\_logs, restrictions \\
Integration & alpaca\_accounts, alpaca\_orders, Partner\_checks, provider\_configs \\
System & configs, cron\_logs, webhooks, notifications \\
\bottomrule
\end{tabular}
\end{table}

\subsection{API Surface}

The 64 routes follow a consistent pattern under \texttt{/v1/}:

\begin{lstlisting}[language=Go,caption={Route registration pattern.}]
func RegisterRoutes(app *base.App) {
    v1 := app.Group("/v1")

    // Trading
    v1.POST("/orders", createOrder)
    v1.GET("/orders/:id", getOrder)
    v1.DELETE("/orders/:id", cancelOrder)
    v1.GET("/orderbook/:security", getOrderBook)

    // Securities
    v1.GET("/securities", listSecurities)
    v1.GET("/securities/:id", getSecurityDetail)
    v1.POST("/securities/:id/transfer", initiateTransfer)

    // Settlement
    v1.POST("/settlements/execute", executeSettlement)
    v1.GET("/settlements/:id/status", getSettlementStatus)

    // Compliance
    v1.POST("/kyc/initiate", initiateKYC)
    v1.GET("/compliance/restrictions/:investor", getRestrictions)
}
\end{lstlisting}

All routes require Hanzo IAM JWT authentication. The \texttt{owner} claim scopes all queries to the authenticated organization.

\subsection{Cron Schedule}

The 19 scheduled tasks maintain system invariants:

\begin{table}[H]
\centering
\caption{Cron job schedule.}
\label{tab:crons}
\begin{tabular}{llp{5.5cm}}
\toprule
\textbf{Job} & \textbf{Schedule} & \textbf{Function} \\
\midrule
settle\_pending & every 30s & Execute pending atomic settlements on Lux EVM \\
match\_orders & every 5s & Run matching engine against open order books \\
sync\_alpaca & every 1m & Synchronize state with Alpaca OmniSub \\
kyc\_poll & every 5m & Poll Partner for KYC/AML status updates \\
nav\_calculate & daily 00:00 UTC & Calculate NAV for fund securities \\
compliance\_scan & every 15m & Check all positions against restriction rules \\
dividend\_distribute & daily 06:00 UTC & Distribute dividends for eligible securities \\
report\_generate & daily 23:00 UTC & Generate regulatory reports (OATS, CAT) \\
mpc\_health & every 1m & Verify MPC signer liveness \\
lockup\_expire & every 1h & Release securities from expired lockup periods \\
accreditation\_check & daily 02:00 UTC & Verify investor accreditation currency \\
market\_data & every 10s & Record market data snapshots \\
fee\_sweep & daily 04:00 UTC & Sweep collected trading fees \\
position\_reconcile & every 30m & Reconcile on-chain state with ATS records \\
audit\_export & weekly Mon 01:00 & Export audit trail for compliance review \\
corporate\_action & every 1h & Process pending corporate actions \\
margin\_call & every 5m & Evaluate margin positions per \cite{lux-credit-protocol-spec} \\
signer\_rotate & monthly 1st 00:00 & Initiate MPC key rotation ceremony \\
backup\_state & every 6h & Backup ATS state to encrypted off-site storage \\
\bottomrule
\end{tabular}
\end{table}

\section{Non-Custodial MPC Custody}\label{sec:mpc}

\subsection{Threshold Signing Architecture}

The central innovation of this ATS is non-custodial operation through MPC threshold signing. The \texttt{luxfi/mpc} library implements a 2-of-3 threshold scheme based on the GG20 protocol for ECDSA and FROST for EdDSA signatures.

Three key shares are distributed:

\begin{enumerate}[nosep]
  \item \textbf{Investor share}: Held on the investor's device, never transmitted to the platform.
  \item \textbf{Platform share}: Held by the ATS in an HSM-backed enclave, used only upon authenticated instruction.
  \item \textbf{Recovery share}: Held by a regulated custodian (e.g., a qualified custodian under the Investment Advisers Act), used only for key recovery.
\end{enumerate}

Any 2 of these 3 shares can produce a valid signature. The platform alone cannot move assets. The investor alone cannot bypass compliance checks. Both must cooperate for a transfer to execute.

\begin{definition}[Non-Custodial ATS]
An ATS is \emph{non-custodial} if for every asset transfer $\tau$, the platform's key share alone is insufficient to produce a valid signature $\sigma$ authorizing $\tau$. Formally, for threshold $t = 2$ and platform share index $i_p$, the signing function $\textsf{Sign}(\{s_{i_p}\}, m) = \bot$ for any message $m$.
\end{definition}

\subsection{Transfer Flow}

A securities transfer proceeds through the following steps:

\begin{algorithm}[H]
\caption{Non-custodial securities transfer.}
\label{alg:transfer}
\begin{algorithmic}[1]
\Require Order $o$, buyer $B$, seller $S$, security token $T$
\Ensure Atomic settlement on Lux EVM
\State ATS verifies $B$ passes KYC/AML and accreditation checks
\State ATS verifies $S$ holds sufficient $T$ in partition $p$
\State ATS verifies transfer satisfies all restriction rules on $T$
\State ATS constructs settlement transaction $\textsf{tx}$
\State ATS requests MPC partial signature from $S$ (investor share)
\State $S$ reviews $\textsf{tx}$ details on device, provides partial signature $\sigma_S$
\State ATS provides platform partial signature $\sigma_P$
\State $\sigma \gets \textsf{Combine}(\sigma_S, \sigma_P)$ \Comment{Valid 2-of-3 threshold signature}
\State ATS submits signed $\textsf{tx}$ to Lux EVM
\State Lux EVM executes atomic transfer: $T$ from $S \to B$, payment from $B \to S$
\State Settlement confirmed in same block (T+0)
\end{algorithmic}
\end{algorithm}

\subsection{Security Properties}

\begin{theorem}[Platform Non-Access]
Under the 2-of-3 threshold scheme, the platform cannot unilaterally transfer any investor asset, even if the platform's HSM is fully compromised.
\end{theorem}

\begin{proof}
The platform holds exactly one key share $s_P$. A valid signature requires $t = 2$ shares. Compromising $s_P$ yields only one share. The adversary must additionally compromise either the investor's device (share $s_I$) or the recovery custodian (share $s_R$). By assumption, these are independent security domains. The platform alone is insufficient.
\end{proof}

\begin{theorem}[Regulatory Compliance]
The MPC architecture preserves the platform's ability to enforce transfer restrictions and comply with regulatory orders, despite being non-custodial.
\end{theorem}

\begin{proof}
The platform's share is required for any 2-of-3 signature involving the platform. The platform will only provide its partial signature when the transfer satisfies all compliance checks (KYC/AML, accreditation, transfer restrictions, court orders). If a regulatory order prohibits a transfer, the platform withholds its share, making the transfer impossible through the normal (investor + platform) signing path. The recovery custodian share exists solely for key recovery under legally defined circumstances and cannot be used for trading.
\end{proof}

\section{Order Matching}\label{sec:matching}

\subsection{Matching Engine}

The matching engine implements price-time priority for continuous order matching:

\begin{lstlisting}[language=Go,caption={Matching engine core.}]
type MatchEngine struct {
    books map[SecurityID]*OrderBook
    mu    sync.RWMutex
}

type OrderBook struct {
    Bids *PriorityQueue // max-heap by price, FIFO by time
    Asks *PriorityQueue // min-heap by price, FIFO by time
}

func (e *MatchEngine) Match(secID SecurityID) []Execution {
    e.mu.Lock()
    defer e.mu.Unlock()

    book := e.books[secID]
    var executions []Execution

    for book.Bids.Len() > 0 && book.Asks.Len() > 0 {
        bestBid := book.Bids.Peek()
        bestAsk := book.Asks.Peek()

        if bestBid.Price.LessThan(bestAsk.Price) {
            break
        }

        qty := min(bestBid.Remaining, bestAsk.Remaining)
        price := bestAsk.Price // price-time priority: aggressor pays

        executions = append(executions, Execution{
            BuyOrder:  bestBid.ID,
            SellOrder: bestAsk.ID,
            Quantity:  qty,
            Price:     price,
            Timestamp: time.Now().UTC(),
        })

        bestBid.Remaining -= qty
        bestAsk.Remaining -= qty

        if bestBid.Remaining == 0 { book.Bids.Pop() }
        if bestAsk.Remaining == 0 { book.Asks.Pop() }
    }
    return executions
}
\end{lstlisting}

\subsection{Order Types}

The ATS supports the following order types, consistent with standard equity market conventions:

\begin{itemize}[nosep]
  \item \textbf{Limit}: Execute at specified price or better.
  \item \textbf{Market}: Execute immediately at best available price.
  \item \textbf{IOC} (Immediate-or-Cancel): Fill what is available, cancel remainder.
  \item \textbf{GTC} (Good-Till-Cancelled): Remain on book until filled or cancelled.
  \item \textbf{GTD} (Good-Till-Date): Remain on book until specified date.
  \item \textbf{RFQ} (Request-for-Quote): Solicit quotes from registered market makers.
\end{itemize}

\section{Settlement}\label{sec:settlement}

\subsection{Atomic T+0 Settlement}

Settlement occurs atomically on Lux EVM. A settlement transaction bundles the securities transfer and the payment transfer into a single atomic operation:

\begin{lstlisting}[language=Solidity,caption={Atomic settlement contract.}]
contract AtomicSettlement {
    function settle(
        address securityToken,
        address paymentToken,
        address seller,
        address buyer,
        uint256 securityAmount,
        uint256 paymentAmount,
        bytes calldata mpcSignature
    ) external {
        require(
            verifyMPCSignature(
                keccak256(abi.encodePacked(
                    securityToken, paymentToken,
                    seller, buyer,
                    securityAmount, paymentAmount
                )),
                mpcSignature
            ),
            "invalid MPC signature"
        );

        IERC1400(securityToken).transferByPartition(
            DEFAULT_PARTITION,
            seller,
            buyer,
            securityAmount,
            ""
        );

        IERC20(paymentToken).transferFrom(
            buyer, seller, paymentAmount
        );

        emit Settled(
            securityToken, seller, buyer,
            securityAmount, paymentAmount,
            block.timestamp
        );
    }
}
\end{lstlisting}

\subsection{Settlement Finality}

Lux consensus provides probabilistic finality within 1--2 seconds. For securities settlement, the ATS waits for deterministic finality (block confirmation) before updating the settlement status. The settlement timeline comparison:

\begin{table}[H]
\centering
\caption{Settlement time comparison.}
\label{tab:settlement}
\begin{tabular}{lll}
\toprule
\textbf{System} & \textbf{Settlement Time} & \textbf{Counterparty Risk} \\
\midrule
Traditional equities (DTCC) & T+1 & Full during settlement period \\
tZERO & T+0 (claimed) & Custodial during settlement \\
Securitize Markets & T+1 to T+3 & Custodial during settlement \\
INX & T+0 to T+2 & Custodial during settlement \\
\textbf{Lux ATS} & \textbf{T+0 atomic} & \textbf{None (non-custodial)} \\
\bottomrule
\end{tabular}
\end{table}

\section{Integration Layer}\label{sec:integration}

\subsection{Alpaca OmniSub}

For traditional equities and hybrid securities, the ATS integrates with Alpaca via OmniSub for brokerage services:

\begin{itemize}[nosep]
  \item Account creation and management via Alpaca Broker API
  \item Order routing for listed equities (NYSE, NASDAQ)
  \item Position synchronization between on-chain and off-chain holdings
  \item Fractional share support for tokenized equity wrappers
\end{itemize}

\subsection{Partner KYC/AML}

Investor identity verification uses Partner for:

\begin{itemize}[nosep]
  \item Document verification (government-issued ID, proof of address)
  \item Biometric authentication (liveness detection, face matching)
  \item AML screening (OFAC, PEP, adverse media)
  \item Accredited investor verification (income, net worth, professional certifications)
  \item Ongoing monitoring (transaction monitoring, periodic re-verification)
\end{itemize}

KYC status is stored in the \texttt{kyc\_checks} collection and linked to the investor's I-Chain DID, enabling portable identity across the Lux ecosystem \cite{lux-transfer-agent}.

\subsection{I-Chain Identity}

Investor identity on the Lux ATS is anchored to the I-Chain decentralized identity system. Each investor has a DID (Decentralized Identifier) that links their KYC/AML verification status, accreditation credentials, and jurisdictional permissions. Transfer restriction checks reference the investor's DID claims rather than centralized database records, enabling:

\begin{itemize}[nosep]
  \item Portable identity across multiple securities platforms on Lux
  \item Zero-knowledge proofs of accreditation without revealing underlying documents
  \item Revocable credentials for expired accreditations or regulatory sanctions
\end{itemize}

\section{Platform Comparison}\label{sec:comparison}

\begin{table}[H]
\centering
\caption{Comparison with existing digital securities platforms.}
\label{tab:comparison}
\begin{tabular}{lcccc}
\toprule
\textbf{Feature} & \textbf{Lux ATS} & \textbf{tZERO} & \textbf{Securitize} & \textbf{INX} \\
\midrule
Non-custodial & Yes & No & No & No \\
MPC threshold signing & 2-of-3 & N/A & N/A & N/A \\
Atomic settlement & T+0 & T+0$^*$ & T+1--3 & T+0--2 \\
Post-quantum ready & Yes & No & No & No \\
On-chain cap table & Yes & Partial & Off-chain & Partial \\
Open-source protocol & Yes & No & No & No \\
Multi-jurisdiction & Yes & US only & US + EU & US + EU \\
ERC-1400 compliant & Yes & Proprietary & ERC-3643 & Proprietary \\
DID-based identity & Yes & No & No & No \\
Integrated order router & Yes & Limited & No & No \\
\bottomrule
\end{tabular}
\end{table}

{\footnotesize $^*$tZERO claims T+0 but requires custodial intermediary during settlement.}

\section{Security Analysis}\label{sec:security}

\subsection{Threat Model}

We consider the following adversaries:

\begin{enumerate}[nosep]
  \item \textbf{External attacker}: Attempts to steal assets by compromising platform infrastructure.
  \item \textbf{Malicious insider}: Platform employee attempts unauthorized transfers.
  \item \textbf{Compromised HSM}: Platform's HSM is fully compromised, exposing the platform key share.
  \item \textbf{Regulatory seizure}: Government orders asset freeze or forfeiture.
  \item \textbf{Quantum adversary}: Future attacker with access to a cryptographically relevant quantum computer.
\end{enumerate}

\subsection{Mitigation}

\begin{table}[H]
\centering
\caption{Threat mitigation through non-custodial MPC.}
\label{tab:threats}
\begin{tabular}{lp{7cm}}
\toprule
\textbf{Threat} & \textbf{Mitigation} \\
\midrule
External attacker & Platform key share alone insufficient; attacker must also compromise investor device \\
Malicious insider & Same as external attacker; no single insider has access to 2 shares \\
Compromised HSM & Investor retains independent share; recovery custodian can facilitate key rotation \\
Regulatory seizure & Platform can freeze by withholding its share; no need to seize keys \\
Quantum adversary & Lux post-quantum precompiles enable migration to PQ-safe threshold schemes \cite{lux-mchain-mpc} \\
\bottomrule
\end{tabular}
\end{table}

\subsection{Operational Security}

\begin{itemize}[nosep]
  \item MPC key shares are generated in a distributed ceremony; no single party ever sees the full private key.
  \item The platform's share is stored in an HSM with FIPS 140-2 Level 3 certification.
  \item Key rotation occurs monthly via the \texttt{signer\_rotate} cron, using proactive secret sharing to rotate shares without changing the public key.
  \item All MPC signing sessions are logged immutably in the \texttt{mpc\_requests} and \texttt{mpc\_signatures} collections.
\end{itemize}

\section{Regulatory Reporting}\label{sec:reporting}

The ATS generates the following regulatory reports:

\begin{itemize}[nosep]
  \item \textbf{OATS} (Order Audit Trail System): All order events with timestamps, required by FINRA.
  \item \textbf{CAT} (Consolidated Audit Trail): Customer and order event information, required by SEC Rule 613.
  \item \textbf{Form ATS-N}: Quarterly operational reports filed with the SEC.
  \item \textbf{FOCUS Reports}: Monthly and quarterly financial reports, required by FINRA Rule 4524.
  \item \textbf{SAR/CTR}: Suspicious Activity Reports and Currency Transaction Reports per FinCEN requirements.
  \item \textbf{Blue Sky}: State-level notice filings for securities offerings.
\end{itemize}

All reports are generated from on-chain data, ensuring that regulatory submissions exactly reflect the executed state of the system. The immutability of on-chain records eliminates the possibility of post-hoc report manipulation.

\section{Performance}\label{sec:performance}

\begin{table}[H]
\centering
\caption{ATS performance characteristics.}
\label{tab:performance}
\begin{tabular}{lr}
\toprule
\textbf{Metric} & \textbf{Value} \\
\midrule
Order submission latency (API) & $< 5$ ms \\
Matching engine throughput & 10,000 orders/sec \\
MPC signing latency (2-of-3) & $< 200$ ms \\
Settlement finality (Lux EVM) & 1--2 seconds \\
End-to-end trade lifecycle & $< 3$ seconds \\
Concurrent order books & 1,000+ \\
Binary size (single Go binary) & $\sim$45 MB \\
Memory footprint (idle) & $\sim$128 MB \\
Memory footprint (peak) & $\sim$512 MB \\
\bottomrule
\end{tabular}
\end{table}

\section{Conclusion}

We have presented the architecture of a non-custodial ATS built on Lux Network that eliminates custodial risk through MPC threshold signing while maintaining full regulatory compliance with SEC Regulation ATS, broker-dealer requirements, and transfer agent obligations. The system descends directly from the Arca securities infrastructure (2018) and represents the culmination of seven years of work on digital securities.

The key architectural decisions---single Go binary on Hanzo Base, 2-of-3 MPC threshold signing via \texttt{luxfi/mpc}, atomic T+0 settlement on Lux EVM, DID-based investor identity---together produce a system that is simultaneously the safest (non-custodial), most decentralized (no single custodian), and most regulatory-compliant (SEC-registered ATS/BD/TA) digital securities platform in operation.

The companion papers on the security token standard \cite{lux-security-token-standard}, smart order routing \cite{lux-smart-order-routing}, and transfer agent architecture \cite{lux-transfer-agent} describe the subsystems in detail. Together, these papers constitute a complete specification for regulatory-compliant, non-custodial securities infrastructure on Lux Network.

\begin{thebibliography}{20}

\bibitem{lux-mchain-mpc}
Z.~Kelling, V.~Seesahai,
``M-Chain: Decentralized Multi-Party Computation Custody with Quantum-Safe Threshold Signatures,''
Lux Partners, 2023.

\bibitem{lux-security-token-standard}
Z.~Kelling,
``Lux Security Token Standard: ERC-1400 on Post-Quantum EVM with MPC Custody,''
Lux Partners, 2025.

\bibitem{lux-smart-order-routing}
Z.~Kelling,
``Smart Order Routing for Digital Securities: Multi-Provider Execution on Lux Network,''
Lux Partners, 2025.

\bibitem{lux-transfer-agent}
Z.~Kelling,
``Blockchain Transfer Agent: On-Chain Cap Table with MPC-Custodied Registry,''
Lux Partners, 2025.

\bibitem{lux-credit-protocol-spec}
Z.~Kelling,
``Lux Credit Protocol Specification,''
Lux Partners, 2023.

\bibitem{sec-reg-ats}
Securities and Exchange Commission,
``Regulation of Exchanges and Alternative Trading Systems,''
Release No.~34-40760, 1998.

\bibitem{sec-rule-15c3-3}
Securities and Exchange Commission,
``Customer Protection---Reserves and Custody of Securities,''
Rule 15c3-3, Securities Exchange Act of 1934.

\bibitem{finra-cat}
Financial Industry Regulatory Authority,
``Consolidated Audit Trail (CAT),''
FINRA Rule 6800 Series.

\bibitem{erc1400}
A.~Dossa, P.~Ruiz, F.~Vogelsteller, S.~Gosselin,
``ERC-1400: Security Token Standard,''
Ethereum Improvement Proposals, 2018.

\bibitem{erc1404}
R.~Bose, J.~Cawley,
``ERC-1404: Simple Restricted Token Standard,''
Ethereum Improvement Proposals, 2018.

\bibitem{gg20}
R.~Gennaro, S.~Goldfeder,
``One Round Threshold ECDSA with Identifiable Abort,''
Cryptology ePrint Archive, 2020.

\bibitem{frost}
C.~Komlo, I.~Goldberg,
``FROST: Flexible Round-Optimized Schnorr Threshold Signatures,''
SAC 2020.

\end{thebibliography}

\end{document}
